% Modernes Inhaltsverzeichnis und Schriftarten für LuaLaTeX
% Nutzt fontspec statt veralteter Font-Befehle

% --- TOC-Paket (tocbasic ist bereits in KOMA enthalten) ---
% tocbasic wird automatisch von KOMA-Script geladen, explizit ist nicht nötig
% \usepackage{tocbasic}  % Bereits integriert

% --- KRITISCH: Alle \fontfamily{phv} ENTFERNT! ---
% Diese funktionieren nicht mit LuaLaTeX + fontspec!
% Die Sans-Serif-Einstellung erfolgt bereits in Schriftart-settings.tex

% --- Moderne Schriftart-Definitionen für Strukturelemente ---
% Da wir bereits Sans-Serif als Standard haben, nur noch Gewichtung/Größe anpassen

% === Hauptstruktur (Kapitel, Sections) ===
\addtokomafont{disposition}{\bfseries}           % Alle Gliederungselemente fett
\addtokomafont{chapter}{\LARGE\bfseries\color{headingcolor}}
\addtokomafont{section}{\Large\bfseries\color{headingcolor}}
\addtokomafont{subsection}{\large\bfseries\color{headingcolor}}
\addtokomafont{subsubsection}{\normalsize\bfseries\color{headingcolor}}
\addtokomafont{paragraph}{\normalsize\bfseries\color{headingcolor}}
\addtokomafont{subparagraph}{\small\bfseries\color{headingcolor}}

% === Titelseite ===
\addtokomafont{title}{\huge\bfseries\color{headingcolor}}
\addtokomafont{author}{\Large\color{secondary}}
\addtokomafont{date}{\large\color{secondary}}

% === Kopf- und Fußzeile ===
\addtokomafont{pageheadfoot}{\small\color{secondary}}
\addtokomafont{pagehead}{\small\color{secondary}}
\addtokomafont{pagefoot}{\footnotesize\color{secondary}}
\addtokomafont{pagenumber}{\small\bfseries}

% === Beschriftungen ===
\addtokomafont{caption}{\small\color{secondary}}
\addtokomafont{captionlabel}{\small\bfseries\color{headingcolor}}

% === Listen und Beschreibungen ===
\addtokomafont{descriptionlabel}{\bfseries\color{headingcolor}}

% === Fußnoten ===
\addtokomafont{footnote}{\footnotesize}
\addtokomafont{footnotelabel}{\footnotesize\bfseries}
\addtokomafont{footnotereference}{\footnotesize}

% === Inhaltsverzeichnis-spezifische Anpassungen ===
\addtokomafont{chapterentry}{\bfseries\color{headingcolor}}
\addtokomafont{chapterentrypagenumber}{\bfseries}

% --- Moderne TOC-Konfiguration ---
\KOMAoptions{
  toc=chapterentrywithdots,      % Punkte auch bei Kapiteln
  toc=flat                       % Flache Hierarchie (modern)
  %listof=chapterentry            % Verzeichnisse wie Kapitel
}

% --- Erweiterte TOC-Anpassungen (vereinfacht) ---
% Abstände und Layout - VEREINFACHT da problematisch

% \RedeclareSectionCommand[
%   tocindent=0pt,                 % Keine Einrückung
%   tocnumwidth=2.5em,             % Breitere Nummern
%   tocentryformat=\bfseries       % Fette Einträge
% ]{chapter}

% \RedeclareSectionCommand[
%   tocindent=1.5em,
%   tocnumwidth=2.3em
% ]{section}

% \RedeclareSectionCommand[
%   tocindent=3.8em,
%   tocnumwidth=3.2em,
%   tocentryformat=\small          % Kleinere Subsections
% ]{subsection}

% --- Moderne Verzeichnis-Titel (polyglossia-kompatibel) ---
% Manuelle Definition da polyglossia statt babel verwendet wird
\renewcommand{\contentsname}{Inhaltsverzeichnis}
\renewcommand{\listfigurename}{Abbildungsverzeichnis}
\renewcommand{\listtablename}{Tabellenverzeichnis}

% Weitere deutsche Bezeichnungen
\renewcommand{\figurename}{Abbildung}
\renewcommand{\tablename}{Tabelle}
\renewcommand{\chaptername}{Kapitel}
\renewcommand{\appendixname}{Anhang}

% --- Custom TOC-Styles (optional) ---
% Für noch moderneres Design - auskommentiert da problematisch
% \newcommand{\moderntocformat}{%
%   \renewcommand*{\@pnumwidth}{2em}     % Seitenzahl-Breite
%   \renewcommand*{\@tocrmarg}{2.5em}    % Rechter Rand
%   \renewcommand*{\@dotsep}{1.5}        % Punktabstand
% }

% Aktiviere moderne Formatierung - DEAKTIVIERT
% \AtBeginDocument{\moderntocformat}

% --- Farb-Integration (vereinfacht) ---
% Links im TOC - DEAKTIVIERT da problematisch
% \makeatletter
% \renewcommand*\l@chapter[2]{%
%   \ifnum \c@tocdepth >\m@ne
%     \addpenalty{-\@highpenalty}%
%     \vskip 1.0em \@plus\p@
%     \setlength\@tempdima{1.5em}%
%     \begingroup
%       \parindent \z@ \rightskip \@pnumwidth
%       \parfillskip -\@pnumwidth
%       \leavevmode \bfseries\color{headingcolor}
%       \advance\leftskip\@tempdima
%       \hskip -\leftskip
%       #1\nobreak\hfil \nobreak\hb@xt@\@pnumwidth{\hss #2}\par
%       \penalty\@highpenalty
%     \endgroup
%   \fi}
% \makeatother

% --- Alternative: Moderne TOC mit tocloft ---
% \usepackage{tocloft}
% \renewcommand{\cftchapfont}{\bfseries\color{headingcolor}}
% \renewcommand{\cftsecfont}{\color{headingcolor}}
% \renewcommand{\cftsubsecfont}{\small\color{secondary}}

% --- Erweiterte Verzeichnisse ---
% Für Code-Listings
\makeatletter
\@ifpackageloaded{listings}{%
  \renewcommand{\lstlistlistingname}{Codeverzeichnis}
  \addtokomafont{lstlistingname}{\bfseries\color{headingcolor}}
}{}
\makeatother

% --- Abkürzungsverzeichnis Integration ---
\makeatletter
\@ifpackageloaded{glossaries}{%
  \renewcommand{\glossaryname}{Glossar}
  \renewcommand{\acronymname}{Abkürzungsverzeichnis}
}{}
\makeatother

% --- Debug-Ausgabe ---
\AtBeginDocument{%
  \typeout{=== TOC-Style konfiguriert ===}
  \typeout{Font-System: fontspec (modern)}
  \typeout{Schrift: Sans-Serif durchgängig}
  \typeout{Farben: Integriert mit Farbschema}
  \typeout{Layout: KOMA + moderne Anpassungen}
  \typeout{=================================}
}